%% ============================================================
%% 通用表格插入模板
%% 变量:
%%   table_caption    — 表格标题
%%   table_label      — 引用标签
%%   table_headers    — 列标题列表 ["指标", "权重", ...]
%%   table_rows       — 行数据 [["x1", "0.32", ...], ...]
%%   table_note       — 表注 (可选)
%% ============================================================

<%- macro render_table(table_caption, table_label, table_headers, table_rows, table_note="") %>
\begin{table}[H]
    \centering
    \caption{<< table_caption >>}
    \label{<< table_label >>}
    \begin{tabular}{<< "c" * table_headers|length >>}
        \toprule
        <% for h in table_headers %><% if not loop.first %> & <% endif %>\textbf{<< h >>}<% endfor %> \\
        \midrule
<% for row in table_rows %>
        <% for cell in row %><% if not loop.first %> & <% endif %><< cell >><% endfor %> \\
<% endfor %>
        \bottomrule
    \end{tabular}
<% if table_note %>
    \begin{tablenotes}
        \small
        \item 注：<< table_note >>
    \end{tablenotes}
<% endif %>
\end{table}
<%- endmacro %>


%% ============================================================
%% 通用长表格插入模板 (跨页)
%% 变量同上
%% ============================================================

<%- macro render_longtable(table_caption, table_label, table_headers, table_rows) %>
\begin{longtable}{<< "c" * table_headers|length >>}
    \caption{<< table_caption >>} \label{<< table_label >>} \\
    \toprule
    <% for h in table_headers %><% if not loop.first %> & <% endif %>\textbf{<< h >>}<% endfor %> \\
    \midrule
    \endfirsthead
    \multicolumn{<< table_headers|length >>}{c}{\tablename\ \thetable\ (续)} \\
    \toprule
    <% for h in table_headers %><% if not loop.first %> & <% endif %>\textbf{<< h >>}<% endfor %> \\
    \midrule
    \endhead
    \bottomrule
    \endfoot
<% for row in table_rows %>
    <% for cell in row %><% if not loop.first %> & <% endif %><< cell >><% endfor %> \\
<% endfor %>
\end{longtable}
<%- endmacro %>


%% ============================================================
%% 通用单图插入模板
%% 变量:
%%   fig_path     — 图片路径
%%   fig_caption  — 图片标题
%%   fig_label    — 引用标签
%%   fig_width    — 宽度 (默认 0.85\textwidth)
%% ============================================================

<%- macro render_figure(fig_path, fig_caption, fig_label, fig_width="0.85\\textwidth") %>
\begin{figure}[H]
    \centering
    \includegraphics[width=<< fig_width >>]{<< fig_path >>}
    \caption{<< fig_caption >>}
    \label{<< fig_label >>}
\end{figure}
<%- endmacro %>


%% ============================================================
%% 双图并排插入模板
%% 变量:
%%   fig1_path, fig1_caption, fig1_label
%%   fig2_path, fig2_caption, fig2_label
%%   overall_caption, overall_label
%% ============================================================

<%- macro render_subfigures(fig1_path, fig1_caption, fig1_label, fig2_path, fig2_caption, fig2_label, overall_caption, overall_label) %>
\begin{figure}[H]
    \centering
    \begin{subfigure}[b]{0.48\textwidth}
        \centering
        \includegraphics[width=\textwidth]{<< fig1_path >>}
        \caption{<< fig1_caption >>}
        \label{<< fig1_label >>}
    \end{subfigure}
    \hfill
    \begin{subfigure}[b]{0.48\textwidth}
        \centering
        \includegraphics[width=\textwidth]{<< fig2_path >>}
        \caption{<< fig2_caption >>}
        \label{<< fig2_label >>}
    \end{subfigure}
    \caption{<< overall_caption >>}
    \label{<< overall_label >>}
\end{figure}
<%- endmacro %>


%% ============================================================
%% 代码块插入模板
%% 变量:
%%   code_content  — 代码文本
%%   code_caption  — 标题
%%   code_label    — 引用标签
%%   code_language — 语言 (默认 Python)
%% ============================================================

<%- macro render_code(code_content, code_caption, code_label, code_language="Python") %>
\begin{lstlisting}[language=<< code_language >>, caption=<< code_caption >>, label=<< code_label >>]
<< code_content >>
\end{lstlisting}
<%- endmacro %>